% !TeX root = ../../main.tex

\section{Introducción}

Una Arquitectura de Big Data establece cómo se recopilan, almacenan,
procesan, protegen y analizan grandes cantidades de datos. Su diseño permite
trabajar con información proveniente de diferentes fuentes y obtener
resultados útiles para apoyar la toma de decisiones.

\section{Definición y propósito}

La Arquitectura de Big Data es el diseño general de los componentes,
tecnologías y procesos utilizados para administrar grandes cantidades de
datos. Su propósito fundamental es procesar información que presenta un gran
volumen, variedad y velocidad.

A diferencia de una arquitectura tradicional, normalmente diseñada para datos
estructurados y cargas relativamente predecibles, una arquitectura de Big
Data puede trabajar con datos estructurados, semiestructurados y no
estructurados. Además, distribuye el almacenamiento y el procesamiento entre
varios equipos, lo que facilita aumentar su capacidad.

\section{Componentes esenciales}

Los componentes principales de una Arquitectura de Big Data son los
siguientes:

\begin{enumerate}
    \item \textbf{Fuentes de datos:} generan la información que será
    procesada. Algunos ejemplos son bases de datos, redes sociales,
    aplicaciones, sensores, archivos y registros de actividad.

    \item \textbf{Ingesta e integración:} recibe los datos y los transporta
    hacia la plataforma. Puede trabajar por lotes o en tiempo real.
    Apache Kafka y Apache NiFi son ejemplos de herramientas de ingesta.

    \item \textbf{Almacenamiento:} conserva los datos de manera distribuida.
    Se pueden utilizar sistemas como HDFS, bases de datos NoSQL, Data Lakes
    y Data Warehouses.

    \item \textbf{Procesamiento:} limpia, transforma y combina los datos.
    Puede realizarse por lotes con Hadoop o mediante procesamiento en tiempo
    real con Spark y Flink.

    \item \textbf{Análisis:} aplica consultas, estadísticas o modelos de
    aprendizaje automático para encontrar patrones y generar conocimiento.

    \item \textbf{Presentación:} muestra los resultados mediante informes,
    gráficas y tableros para facilitar su interpretación.

    \item \textbf{Gobernanza y seguridad:} establece reglas sobre la calidad,
    propiedad, acceso, protección y conservación de los datos.
\end{enumerate}

\section{Arquitectura Lambda}

La Arquitectura Lambda combina el procesamiento por lotes con el
procesamiento en tiempo real. Su objetivo es ofrecer resultados rápidos sin
perder la precisión que se obtiene al analizar todos los datos históricos.

Está formada por tres capas:

\begin{itemize}
    \item \textbf{Capa Batch:} conserva el conjunto histórico de datos y lo
    procesa periódicamente. Genera resultados completos y precisos, aunque
    requiere más tiempo.

    \item \textbf{Capa Speed:} procesa los datos recientes conforme llegan.
    Su función es producir resultados con poca demora mientras la capa Batch
    termina su procesamiento.

    \item \textbf{Capa Serving:} almacena y presenta los resultados generados
    por las otras capas para que puedan ser consultados por aplicaciones y
    usuarios.
\end{itemize}

Un caso de uso es la detección de fraude bancario. La capa Speed puede revisar
una transacción en el momento en que ocurre, mientras que la capa Batch
analiza el historial completo del cliente. La capa Serving permite consultar
ambos resultados para decidir si la transacción parece sospechosa.

\section{Arquitectura Lambda y Arquitectura Kappa}

La Arquitectura Kappa utiliza una sola ruta de procesamiento basada en flujos
de eventos. Los datos históricos pueden volver a procesarse reproduciendo los
eventos almacenados. Esto evita mantener una ruta Batch y otra de tiempo real.

\begin{center}
    \renewcommand{\arraystretch}{1.4}
    \small

    \begin{tabularx}{\textwidth}{|p{2.7cm}|X|X|}
        \hline
        \textbf{Aspecto} &
        \textbf{Arquitectura Lambda} &
        \textbf{Arquitectura Kappa} \\
        \hline

        Procesamiento &
        Utiliza procesamiento por lotes y en tiempo real. &
        Utiliza una sola ruta de procesamiento de eventos. \\
        \hline

        Capas &
        Tiene las capas Batch, Speed y Serving. &
        Se concentra en la ingesta, procesamiento y consulta de flujos. \\
        \hline

        Complejidad &
        Es más compleja porque mantiene dos formas de procesamiento. &
        Es más sencilla porque mantiene una sola forma de procesamiento. \\
        \hline

        Datos históricos &
        Se vuelven a calcular mediante la capa Batch. &
        Se procesan nuevamente reproduciendo el flujo de eventos. \\
        \hline

        Uso recomendado &
        Sistemas que necesitan combinar resultados históricos precisos con
        información reciente. &
        Sistemas donde predominan los eventos y las respuestas en tiempo
        real. \\
        \hline
    \end{tabularx}
\end{center}

Lambda es adecuada cuando se requieren análisis históricos completos y
resultados inmediatos. Kappa es conveniente en aplicaciones orientadas a
eventos, como monitoreo de sensores, actividad de usuarios o transacciones
en tiempo real.

\section{Arquitectura de Data Lake}

Un Data Lake es un repositorio que almacena grandes cantidades de datos en su
formato original. Puede contener datos estructurados, como tablas;
semiestructurados, como archivos JSON o XML; y no estructurados, como
imágenes, videos, audios y documentos.

Una característica importante es el esquema en lectura. Esto significa que
los datos se organizan cuando van a utilizarse, en lugar de definir toda su
estructura antes de almacenarlos.

\subsection{Ventajas}

\begin{itemize}
    \item Permite almacenar diferentes tipos de datos.
    \item Puede crecer conforme aumenta la cantidad de información.
    \item Conserva los datos originales para futuros análisis.
    \item Facilita proyectos de análisis y aprendizaje automático.
    \item Puede utilizar almacenamiento distribuido de bajo costo.
\end{itemize}

\subsection{Desafíos}

\begin{itemize}
    \item Mantener la calidad y organización de los datos.
    \item Crear un catálogo para localizar la información.
    \item Controlar el acceso a datos sensibles.
    \item Evitar que el repositorio se convierta en un conjunto desordenado
    conocido como \textit{data swamp}.
\end{itemize}

\section{Herramientas y técnicas de Big Data}

Las herramientas y técnicas pueden clasificarse en cuatro categorías
principales:

\begin{enumerate}
    \item \textbf{Procesamiento masivamente paralelo:} divide una tarea entre
    varios procesadores para ejecutarla al mismo tiempo. Algunos ejemplos son
    Teradata, Greenplum y Amazon Redshift.

    \item \textbf{Bases de datos NoSQL:} almacenan datos que no siempre
    requieren una estructura relacional fija. Algunos ejemplos son MongoDB,
    Apache Cassandra y Apache HBase.

    \item \textbf{Almacenamiento y procesamiento distribuido:} reparte los
    datos y el trabajo entre varios nodos. Algunos ejemplos son Apache Hadoop,
    HDFS, Apache Spark y Apache Flink.

    \item \textbf{Computación en la nube:} proporciona almacenamiento y
    procesamiento que puede aumentar o disminuir según la demanda. Algunos
    ejemplos son Amazon Web Services, Microsoft Azure y Google Cloud.
\end{enumerate}

\section{Beneficios}

Tres beneficios importantes de implementar una Arquitectura de Big Data son:

\begin{enumerate}
    \item \textbf{Escalabilidad:} permite agregar recursos cuando aumenta la
    cantidad de datos o el trabajo de procesamiento.

    \item \textbf{Mayor velocidad de procesamiento:} distribuye las tareas
    entre varios nodos, lo que permite analizar grandes cantidades de
    información en menos tiempo.

    \item \textbf{Mejores decisiones:} ayuda a identificar tendencias,
    riesgos y oportunidades mediante información obtenida de los datos.
\end{enumerate}

También ofrece flexibilidad para combinar diferentes herramientas y puede
mejorar la disponibilidad mediante la distribución y réplica de los datos.

\section{Desafíos de implementación}

Tres desafíos comunes son:

\begin{enumerate}
    \item \textbf{Complejidad técnica:} integrar y mantener diferentes
    herramientas requiere personal capacitado y una planeación adecuada.

    \item \textbf{Calidad de los datos:} la información puede contener
    errores, duplicados o formatos distintos. Es necesario establecer
    procesos de limpieza y validación.

    \item \textbf{Seguridad y privacidad:} almacenar grandes cantidades de
    información aumenta el riesgo de accesos no autorizados y filtraciones.
    Por ello deben utilizarse controles de acceso, cifrado y monitoreo.
\end{enumerate}

\section{Casos de uso}

La Arquitectura de Big Data se utiliza para transformar grandes cantidades de
información en resultados útiles. Algunos casos de uso importantes son:

\begin{itemize}
    \item \textbf{Detección de fraude:} analiza transacciones en tiempo real
    para identificar comportamientos sospechosos.

    \item \textbf{Sistemas de recomendación:} estudia el comportamiento de
    los usuarios para sugerir productos, películas o servicios.

    \item \textbf{Mantenimiento predictivo:} procesa información de sensores
    para detectar posibles fallas en máquinas y equipos.

    \item \textbf{Análisis comercial:} combina ventas, inventarios y
    comportamiento de clientes para mejorar las decisiones de una empresa.
\end{itemize}

\section{Gobernanza, privacidad y seguridad}

La gobernanza de datos define quién es responsable de la información, quién
puede utilizarla y bajo qué condiciones. También establece normas para
mantener la calidad, clasificación, trazabilidad y conservación de los datos.

Las medidas de protección esenciales incluyen:

\begin{itemize}
    \item Cifrado de los datos almacenados y transmitidos.
    \item Autenticación de usuarios y control de acceso por roles.
    \item Aplicación del principio de menor privilegio.
    \item Enmascaramiento o anonimización de información personal.
    \item Registro y auditoría de las actividades realizadas.
    \item Monitoreo para detectar accesos o comportamientos sospechosos.
    \item Copias de seguridad y planes de recuperación.
    \item Políticas para conservar y eliminar datos cuando corresponda.
\end{itemize}

Estas medidas permiten reducir riesgos, proteger la privacidad y cumplir las
normas aplicables al uso de la información.

\section{Conclusión}

Una Arquitectura de Big Data permite administrar datos que superan las
capacidades de muchas soluciones tradicionales. Sus componentes se encargan
de recibir, almacenar, procesar, analizar y proteger la información.
Arquitecturas como Lambda, Kappa y Data Lake ofrecen alternativas para
diferentes necesidades. La elección debe considerar el tipo de datos, la
velocidad requerida, el costo, la seguridad y los objetivos de la
organización.